\documentclass[11pt,aspectratio=169]{beamer}
\ifdefined\pdfmapfile\pdfmapfile{+cm-super-t1.map}\fi
% Shared preamble for MA3521 lectures.
\usepackage{beamerthemeMTU}
\usepackage{amsmath,amssymb,cancel}
\course{MA3521 --- Introduction to ODEs}
\author{Allan Struthers}
\institute{Math Sciences}
% Reusable three-tank network. Load once in the preamble.
% \ThreeTankDiagram{capacity1=100,volume1=60,concentration1=2,...}
% Volumes: L. Concentrations: g/L. Flow rates: L/min.
\usetikzlibrary{arrows.meta}
\pgfkeys{/ode tanks/.is family,/ode tanks,
  capacity1/.store in=\ODEcapone, capacity2/.store in=\ODEcaptwo, capacity3/.store in=\ODEcapthree,
  volume1/.store in=\ODEvolone, volume2/.store in=\ODEvoltwo, volume3/.store in=\ODEvolthree,
  concentration1/.store in=\ODEcone, concentration2/.store in=\ODEctwo, concentration3/.store in=\ODEcthree,
  inflow/.store in=\ODEqin, inflowconcentration/.store in=\ODEcin,
  flow12/.store in=\ODEqonetwo, flow21/.store in=\ODEqtwoone,
  flow23/.store in=\ODEqtwothree, flow32/.store in=\ODEqthreetwo,
  outflow/.store in=\ODEqout}

% Shared compact labels used by both lectures and worksheets.
\newcommand{\ODETankHeading}[2]{%
  \node[font=\bfseries] at (0,3.12) {$T_{#1}$ (#2 L)};}
\newcommand{\ODETankLabels}[4]{%
  \pgfmathtruncatemacro{\InitialSalt}{#3*#4}
  \node[anchor=north,inner sep=1pt,font=\fontsize{7}{8}\selectfont] at (0,-.20)
  {$v_{#1}(0)=#3\,\mathrm{L},\;s_{#1}(0)=\InitialSalt\,\mathrm{g}$};}
% One cylindrical tank: index, horizontal position, capacity, initial volume,
% initial concentration. The liquid height follows initial volume/capacity.
\newcommand{\ODETank}[5]{%
  \begin{scope}[shift={(#2,0)}]
    \pgfmathsetmacro{\ODElevel}{2.6*(#4)/(#3)}
    \path[fill=mtuGreenPale] (-.9,0) arc[start angle=180,end angle=360,x radius=.9,y radius=.17]
      -- (.9,\ODElevel) arc[start angle=0,end angle=180,x radius=.9,y radius=.17] -- cycle;
    \draw[mtuGreenLight,fill=mtuGreenLight!20] (0,\ODElevel) ellipse[x radius=.9,y radius=.17];
    \draw[thick] (-.9,2.6)--(-.9,0) arc[start angle=180,end angle=360,x radius=.9,y radius=.17] -- (.9,2.6);
    \draw[thick] (0,2.6) ellipse[x radius=.9,y radius=.17];
    \ODETankHeading{#1}{#3}
    \ODETankLabels{#1}{#3}{#4}{#5}
    \coordinate (tank#1wh) at (-.9,1.95);
    \coordinate (tank#1eh) at (.9,1.95);
    \coordinate (tank#1wl) at (-.9,.7);
    \coordinate (tank#1el) at (.9,.7);
  \end{scope}}

\newcommand{\ThreeTankDiagram}[1]{%
  \begingroup
  \pgfkeys{/ode tanks/.cd,
    capacity1=100,capacity2=120,capacity3=100,
    volume1=60,volume2=80,volume3=50,
    concentration1=2,concentration2=1,concentration3=0,
    inflow=2,inflowconcentration=3,flow12=5,flow21=3,flow23=4,flow32=2,outflow=2,#1}
  \begin{tikzpicture}[font=\small,
    pipe/.style={-{Stealth[length=2.5mm,width=1.7mm]},thick,mtuGreen},
    rate/.style={above,align=center,text=black,fill=white,inner sep=3pt}]
    \ODETank{1}{0}{\ODEcapone}{\ODEvolone}{\ODEcone}
    \ODETank{2}{4.6}{\ODEcaptwo}{\ODEvoltwo}{\ODEctwo}
    \ODETank{3}{9.2}{\ODEcapthree}{\ODEvolthree}{\ODEcthree}
    \draw[pipe] (-3.7,1.95)--(tank1wh) node[midway,rate] {$\begin{aligned}q_{\mathrm{in}}&=\ODEqin\ \mathrm{L/min}\\c_{\mathrm{in}}&=\ODEcin\ \mathrm{g/L}\end{aligned}$};
    \draw[pipe] (tank1eh)--(tank2wh) node[midway,rate] {$q_{12}=\ODEqonetwo\ \mathrm{L/min}$};
    \draw[pipe] (tank2wl)--(tank1el) node[midway,rate] {$q_{21}=\ODEqtwoone\ \mathrm{L/min}$};
    \draw[pipe] (tank2eh)--(tank3wh) node[midway,rate] {$q_{23}=\ODEqtwothree\ \mathrm{L/min}$};
    \draw[pipe] (tank3wl)--(tank2el) node[midway,rate] {$q_{32}=\ODEqthreetwo\ \mathrm{L/min}$};
    \draw[pipe] (tank3el)--(12.9,.7) node[midway,rate] {$q_{\mathrm{out}}=\ODEqout\ \mathrm{L/min}$};
  \end{tikzpicture}%
  \endgroup}

% Reusable shear-building sketch. Displacements are absolute, measured from rest.
% Numeric data: masses in kg, stiffnesses in N/m, damping in N s/m.
% Drawing offsets are independent of physical units.
% basemotion accepts a math expression, for example basemotion={\sin(t)}.
\usetikzlibrary{arrows.meta,calc}
\pgfkeys{/ode building/.is family,/ode building,
  baseoffset/.store in=\ODEbase, basemotion/.store in=\ODEbasemotion,
  offset1/.store in=\ODEbone,offset2/.store in=\ODEbtwo,offset3/.store in=\ODEbthree,
  mass1/.store in=\ODEbmone,mass2/.store in=\ODEbmtwo,mass3/.store in=\ODEbmthree,
  stiffness1/.store in=\ODEbkone,stiffness2/.store in=\ODEbktwo,stiffness3/.store in=\ODEbkthree,
  damping1/.store in=\ODEbcone,damping2/.store in=\ODEbctwo,damping3/.store in=\ODEbcthree}

% Dashpot between two numeric points. The piston follows the line of action.
\newcommand{\ODEBuildingDamper}[4]{%
  \pgfmathsetmacro{\ODEangle}{atan2(#4-#2,#3-#1)}
  \pgfmathsetmacro{\ODElength}{sqrt((#3-#1)^2+(#4-#2)^2)}
  \begin{scope}[shift={(#1,#2)},rotate=\ODEangle]
    \draw[thick] (0,0)--({.38*\ODElength},0);
    \draw[thick,fill=white] ({.38*\ODElength},-.075) rectangle ({.72*\ODElength},.075);
    \draw[thick] ({.55*\ODElength},-.065)--({.55*\ODElength},.065);
    \draw[thick] ({.55*\ODElength},0)--(\ODElength,0);
  \end{scope}}

\newcommand{\ThreeStoryBuilding}[1]{%
  \begingroup
  \pgfkeys{/ode building/.cd,
    basemotion={\sin(t)},baseoffset=.30,offset1=.65,offset2=1.05,offset3=1.45,
    mass1=4,mass2=3,mass3=2,
    stiffness1=4,stiffness2=3,stiffness3=2,
    damping1=2,damping2=1,damping3=1,#1}
  \begin{tikzpicture}[font=\normalsize,
    motion/.style={-{Stealth[length=2mm]},thick,mtuGreen},
    ref/.style={gray!55,dashed,thin},
    guide/.style={gray!50,densely dotted,thin},
    tag/.style={fill=white,inner sep=1.5pt}]
    % Displacement arrows extend from the outside datum to the left floor edges.
    % The datum is offset for legibility. Arrows are schematic, not to scale.
    % Shape offsets control the cubic column geometry independently.
    \draw[ref] (-3.15,-.05)--(-3.15,4.37);
    \node[anchor=south,font=\small] at (1.5,4.85) {Displacements: $\vect{y}(t)=(y_1(t),y_2(t),y_3(t))$ (m),\quad $t$ in seconds.};
    \node[gray,anchor=north,font=\small] at (-3.15,-.1) {$0$};
    % Light rest-position outline.

    \foreach \i in {1,2,3}{\draw[ref] (-1.2,{1.35*\i})--(1.2,{1.35*\i});}
    % Slabs end flush with the columns. No overhanging edges.
    \foreach \i/\offset/\mass in {1/\ODEbone/\ODEbmone,2/\ODEbtwo/\ODEbmtwo,3/\ODEbthree/\ODEbmthree}{
      \draw[thick,fill=gray!25] ({\offset-1.2},{1.35*\i-.055}) rectangle ({\offset+1.2},{1.35*\i+.055});


      \draw[motion] (-3.15,{1.35*\i})--({\offset-1.2},{1.35*\i});
      \node[above,mtuGreen] at (-2.6,{1.35*\i}) {$y_{\i}(t)$};

    }
    % Each outside edge is ONE continuous piecewise cubic path. Vertical
    % tangents and matching derivatives at the floor junctions give C1 joins.
    \foreach \side in {-1.2,1.2}{
      \draw[line width=1.7pt] ({\ODEbase+\side},0)
      .. controls ({\ODEbase+\side},.45) and ({\ODEbone+\side},.90) .. ({\ODEbone+\side},1.35)
      .. controls ({\ODEbone+\side},1.80) and ({\ODEbtwo+\side},2.25) .. ({\ODEbtwo+\side},2.70)
      .. controls ({\ODEbtwo+\side},3.15) and ({\ODEbthree+\side},3.60) .. ({\ODEbthree+\side},4.05);
    }
    \foreach \i/\lo/\hi/\stiff/\damp/\mass in {
      1/\ODEbase/\ODEbone/\ODEbkone/\ODEbcone/\ODEbmone,
      2/\ODEbone/\ODEbtwo/\ODEbktwo/\ODEbctwo/\ODEbmtwo,
      3/\ODEbtwo/\ODEbthree/\ODEbkthree/\ODEbcthree/\ODEbmthree}{
      \pgfmathsetmacro{\ODElowy}{1.35*(\i-1)}
      \pgfmathsetmacro{\ODEhighy}{1.35*\i}
      \pgfmathsetmacro{\ODEmid}{(\lo+\hi)/2}
      \node[tag] at ({\ODEmid-.87},{\ODElowy+.67}) {$k_{\i}$};
      \ODEBuildingDamper{\lo+.2}{\ODElowy+.055}{\hi+.65}{\ODEhighy-.055}
      \node[anchor=west,tag] at ({\ODEmid+.52},{\ODElowy+.66}) {$c_{\i}$};
      \node[anchor=west,font=\small,draw=mtuGreen,line width=.6pt,inner xsep=4pt,inner ysep=1pt,minimum width=2.7cm] at (3.35,{\ODEhighy-.32+.2*(\i-2)}) {\setlength{\jot}{0pt}$\begin{aligned}
        m_{\i}&=\mass\,\mathrm{kg}\\k_{\i}&=\stiff\,\mathrm{N/m}\\c_{\i}&=\damp\,\mathrm{N\,s/m}
      \end{aligned}$};
    }
    % Foundation and rollers. The baseoffset uses the same outside strip.
    \draw[thick,fill=gray!40] ({\ODEbase-1.5},-.16) rectangle ({\ODEbase+1.5},0);
    \foreach \wheel in {-.95,.95}{\draw[thick,fill=white] ({\ODEbase+\wheel},-.31) circle[radius=.15];}
    \draw[thick] (-1.8,-.47)--(2.4,-.47);
    \foreach \ground in {-1.7,-1.45,...,2.3}{\draw[gray,thin] (\ground,-.47)--({\ground-.12},-.6);}

    \draw[motion] (-3.15,-.02)--({\ODEbase-1.5},-.02);
    \node[above,mtuGreen] at (-2.6,-.02) {$y_0(t)$};
    \node[anchor=west,align=left,font=\small] at (3.35,-.13) {$y_0(t)=\ODEbasemotion\;\mathrm{m}$};
  \end{tikzpicture}%
  \endgroup}

\usepackage{booktabs}
\title{Great Lakes}
\date{}
\lecturelabel{Great Lakes}

\begin{document}

\begin{frame}{Great Lakes}
\centering
% Schematic placement only. Lake shapes and distances are not to scale.
\begin{tikzpicture}[
  lake/.style={draw=mtuGreen,fill=mtuGreenPale,rounded corners=5pt,
    minimum height=0.42cm,inner xsep=6pt,font=\footnotesize},
  connection/.style={->,thick,draw=mtuGreen}]
\node[lake] (S) at (-2.7,0.85) {S};
\node[lake] (M) at (-1.7,-0.10) {M};
\node[lake] (H) at (0,0.40) {H};
\node[lake] (E) at (1.8,-0.20) {E};
\node[lake] (O) at (3.5,0.35) {O};
\draw[connection] (S.east) -- (H.north west);
\draw[<->,thick,draw=mtuGreen] (M.east) -- (H.south west);
\draw[connection] (H.south east) -- (E.west);
\draw[connection] (E.east) -- (O.south west);
\end{tikzpicture}

\par\smallskip
{\footnotesize S = Superior\quad M = Michigan\quad H = Huron\quad E = Erie\quad O = Ontario\par}

\smallskip
\footnotesize
\setlength{\tabcolsep}{3pt}
\renewcommand{\arraystretch}{1.35}
\begin{tabular}{lrrrrlll}
\toprule
Lake & \shortstack{Volume\\$10^{12}\,\mathrm{m^3}$}
 & \shortstack{Precip.\\$10^3\,\mathrm{m^3/s}$}
 & \shortstack{Watershed\\$10^3\,\mathrm{m^3/s}$}
 & \shortstack{Evap.\\$10^3\,\mathrm{m^3/s}$}
 & \shortstack{Lake inflow\\$10^3\,\mathrm{m^3/s}$}
 & \shortstack{Lake outflow\\$10^3\,\mathrm{m^3/s}$}
 & \shortstack{Other outflow\\$10^3\,\mathrm{m^3/s}$} \\
\midrule
S & 12.10 & $+2.00$ & $+1.50$ & $-1.40$ & --- & $-2.10$ H & --- \\
M & 4.92 & $+1.50$ & $+1.29$ & $-1.10$ & $+6.00$ H & $-7.60$ H & $-0.09$ Chicago \\
H & 3.54 & $+1.60$ & $+1.40$ & $-1.10$
 & \shortstack[l]{$+2.10$ S\\$+7.60$ M}
 & \shortstack[l]{$-6.00$ M\\$-5.60$ E} & --- \\
E & 0.484 & $+0.70$ & $+0.70$ & $-0.70$ & $+5.60$ H & $-6.30$ O & --- \\
O & 1.64 & $+0.50$ & $+1.00$ & $-0.40$ & $+6.30$ E & --- & $-7.40$ St. Lawrence \\
\bottomrule
\end{tabular}

\end{frame}

\begin{frame}{Volume balance}
\small
With $t$ in seconds and each $v_i(t)$ in $\mathrm{m^3}$
\[
\begin{aligned}
v_s'(t) &= (+2.00+1.50-1.40-2.10)\times10^3=0\\
v_m'(t) &= (+1.50+1.29-1.10+6.00-7.60-0.09)\times10^3=0\\
v_h'(t) &= (+1.60+1.40-1.10+2.10+7.60-6.00-5.60)\times10^3=0\\
v_e'(t) &= (+0.70+0.70-0.70+5.60-6.30)\times10^3=0\\
v_o'(t) &= (+0.50+1.00-0.40+6.30-7.40)\times10^3=0
\end{aligned}
\]
Lake volumes are constant.
\[
\begin{aligned}
v_s&=12.10\times10^{12}\,\mathrm{m^3}\quad
&v_m&=4.92\times10^{12}\,\mathrm{m^3}\quad
&v_h&=3.54\times10^{12}\,\mathrm{m^3}\\
v_e&=0.484\times10^{12}\,\mathrm{m^3}\quad
&v_o&=1.64\times10^{12}\,\mathrm{m^3}
\end{aligned}
\]
\end{frame}

\begin{frame}{Pollutant balance in well-mixed lakes}
\small
\begin{itemize}
\item Lakes are well-mixed tanks with constant volume $v$ in $\mathrm{m^3}$.
\item Pollutants do not evaporate. There is no external pollutant supply. Amounts change only through water transfers.
\item Pollutant amount is $s(t)$ in $\mathrm{kg}$ and concentration is $s(t)/v$ in $\mathrm{kg/m^3}$. Initial concentration $c_0$ gives $s(0)=vc_0$.
\item An inflow $q$ in $\mathrm{m^3/s}$ carries its source lake's concentration $c$. Total water outflow $q_{\mathrm{out}}$ excludes evaporation and carries this lake's concentration $s(t)/v$.
\end{itemize}

\medskip
Flow times concentration gives pollutant rate
\(
(\cancel{\mathrm{m^3}}/\mathrm{s})(\mathrm{kg}/\cancel{\mathrm{m^3}})
=\mathrm{kg}/\mathrm{s}
\).
\[
s'(t)=\underbrace{\sum_{\text{incoming flows}}qc}_{\text{pollutant inflow rate}}
-\underbrace{q_{\mathrm{out}}s(t)/v}_{\text{pollutant outflow rate}}
\]
\end{frame}

\begin{frame}{Building the five-lake matrix}
\small
Use $\vect{s}=(s_s,s_m,s_h,s_e,s_o)$ with amounts in $\mathrm{kg}$ and time in seconds.
Michigan receives pollutant from Huron and loses it to Huron and Chicago
\[
s_m'=10^{-9}\bigl(6.00s_h/3.54-(7.60+0.09)s_m/4.92\bigr)
\]
Each row is a receiving lake's balance. Each column uses its source lake's volume.
\[
\boxed{\vect{s}'=A\vect{s}}
\qquad
A=10^{-9}
\begin{pmatrix}
-2.10/12.10 & 0 & 0 & 0 & 0\\
0 & -7.69/4.92 & 6.00/3.54 & 0 & 0\\
2.10/12.10 & 7.60/4.92 & -11.60/3.54 & 0 & 0\\
0 & 0 & 5.60/3.54 & -6.30/0.484 & 0\\
0 & 0 & 0 & 6.30/0.484 & -7.40/1.64
\end{pmatrix}
\]
The constant matrix $A$ has entries in $\mathrm{s^{-1}}$ since
\((10^3\,\mathrm{m^3/s})/(10^{12}\,\mathrm{m^3})=10^{-9}\,\mathrm{s^{-1}}\).
\end{frame}

\begin{frame}{Estimation notes}
\begin{itemize}
\item Historical starting values are principally 1953--2010 averages.
\item Rounded teaching flows have been adjusted so every lake has zero net water flow. Transfers match at both ends.
\item Michigan--Huron precipitation and evaporation were initially split by surface area, and runoff by land drainage area.
\item Mackinac directional flows are assumptions. The final pair is $7600$ toward Huron and $6000$ toward Michigan, in $\mathrm{m^3/s}$.
\item Superior watershed inflow includes Long Lac and Ogoki. Huron includes Georgian Bay.
\item The Huron--Erie route is one transfer with no Lake St. Clair tank. Erie--Ontario combines Niagara and Welland.
\item Volumes are reference-level storage estimates.
\end{itemize}
\end{frame}

\begin{frame}{Mackinac exchange estimate}
The initial directional estimates used an assumed oscillating transport
\[
Q(t)=1100+20000\sin(2\pi t/T)\qquad [Q]=\mathrm{m^3/s}
\]
Positive transport is toward Huron. Averaging over a complete cycle gives
\[
\begin{aligned}
\left\langle\max(Q(t),0)\right\rangle &\approx 6930\ \mathrm{m^3/s}\\
\left\langle\max(-Q(t),0)\right\rangle &\approx 5830\ \mathrm{m^3/s}
\end{aligned}
\]
The final teaching values were adjusted to $7600$ and $6000\ \mathrm{m^3/s}$ to balance the five-lake model. These are assumed directional averages, not measured long-term averages.
\end{frame}

\begin{frame}{References}
\begin{enumerate}
\item University of Michigan, Graham Sustainability Institute (2015).\newline
\href{https://graham.umich.edu/media/pubs/GreatLakesWaterBudget.pdf}{\emph{Great Lakes Water Budget}, Figure 1}.\newline
Historical water-budget components and connecting flows from NOAA-GLERL data.
\item U.S. Environmental Protection Agency.\newline
\href{https://www.epa.gov/greatlakes/physical-features-great-lakes}{\emph{Physical Features of the Great Lakes}}.\newline
Surface and land drainage areas for the initial Michigan--Huron split.
\item NOAA Technical Memorandum GLERL-137, Table 3.\newline
\href{https://repository.library.noaa.gov/view/noaa/11063/noaa_11063_DS1.pdf}{\emph{Modified Great Lakes Hydrology Modeling System for Considering Simple Extreme Climates}}.\newline
Coordinated lake volumes, including Georgian Bay within Huron.
\end{enumerate}
\end{frame}

\begin{frame}{References for Mackinac exchange}
\begin{enumerate}
\setcounter{enumi}{3}
\item Anderson and Schwab (2013).\newline
\href{https://doi.org/10.1016/j.jglr.2013.09.001}{\emph{Predicting the oscillating bi-directional exchange flow in the Straits of Mackinac}}.\newline
Net transport near $1100\ \mathrm{m^3/s}$ and reversing exchange.
\item U.S. Environmental Protection Agency.\newline
\href{https://nepis.epa.gov/Exe/ZyPURL.cgi?Dockey=9101AX9F.TXT}{\emph{Biological, Chemical and Physical Relationships in the Straits of Mackinac}}.\newline
Observed transport magnitudes motivating the assumed oscillation amplitude.
\item Anderson and Schwab (2017).\newline
\href{https://doi.org/10.1002/2016JC012255}{\emph{Meteorological influence on summertime baroclinic exchange in the Straits of Mackinac}}.\newline
Simultaneous opposing flows during summer stratification.
\end{enumerate}
\end{frame}

\end{document}
